把逐边失谐放回全相互作用绘景后，可直接估计位移模型遗漏的振幅。令 $\tau=t-t_i$，真实第 $\ell$ 条吸收边的速率为
\[
 v_\ell(t)=-\ii\mathcal G_\ell(t)\ee^{\ii\Delta_\ell\tau},\qquad
 \Delta_\ell=\delta_{\ell+1}-\delta_\ell,
\]
均匀模型以 $v_0(t)=\upsilon(t)$ 代替所有 $v_\ell$。这里只比较逐边变化；物理格边界和被舍环境模式的误差需另外加入。设 $P_W$ 投影到选定电子工作窗，$\mathcal E_W$ 包含与该窗中至少一个格点相接的交换边。由 $|\ee^{\ii x}-1|\le|x|$，每条相关边的速率差满足
\begin{equation}
 \max_{\ell\in\mathcal E_W}|v_\ell(t)-v_0(t)|
 \le a_W(t)\equiv
 \max_{\ell\in\mathcal E_W}|\mathcal G_\ell(t)-\mathcal G_0(t)|
 +|\mathcal G_0(t)|\tau
  \max_{\ell\in\mathcal E_W}|\Delta_\ell-\Delta_0|.
 \label{eq:feqo-quantum-edge-rate-error}
\end{equation}
采用绝对速率差，便不会在脉冲零点或共同矩阵元为零时除以零。

\begin{proposition}[逐边变化与玻色数矩的传播误差]
令 $U$ 为真实逐边交换传播子，$U_0$ 为共同速率的双边位移模型传播子，二者存在于同一联合空间中。对归一化入射态，记 $\Psi_0(t)=U_0(t,t_i)\Psi_i$，$R(t)=H(t)-H_0(t)$。若相应 Duhamel 积分与数矩有限，则
\begin{equation}
\begin{aligned}
 \|(U-U_0)\Psi_i\|
 &\le2\int_{t_i}^{t_f}a_W(t)
       \sqrt{\langle N_{\rm ph}\rangle_{0,t}+1}\,\dd t
       +\epsilon_{W,\rm out},\\
 \epsilon_{W,\rm out}
 &\equiv\frac1\hbar\int_{t_i}^{t_f}
       \|R(t)(I-P_W)\Psi_0(t)\|\dd t.
\end{aligned}
\label{eq:feqo-quantum-ladder-propagation-error}
\end{equation}
这里数矩属于参考演化的实际光子态，不限于初始 Fock 数；交换中的光子产生可以使它增大。
\end{proposition}
\begin{proof}
在全相互作用绘景中把吸收速率差写成加权电子平移
\[
 D_e(t)=\sum_\ell[v_\ell(t)-v_0(t)]
                 |\ell+1\rangle\langle\ell|,
 \qquad R=\ii\hbar(D_ea-D_e^\dagger a^\dagger).
\]
与工作窗相接的各边系数由 $a_W$ 控制，所以
$\|D_eP_W\|,\|D_e^\dagger P_W\|\le a_W$。电子投影与光子算符对易，于是
\begin{align*}
 \frac1\hbar\|RP_W\Psi_0\|
 &\le a_W\bigl(\|aP_W\Psi_0\|
                         +\|a^\dagger P_W\Psi_0\|\bigr)\\
 &\le a_W\bigl(\sqrt{\langle N_{\rm ph}\rangle_{0,t}}
                   +\sqrt{\langle N_{\rm ph}\rangle_{0,t}+1}\bigr)\\
 &\le2a_W\sqrt{\langle N_{\rm ph}\rangle_{0,t}+1}.
\end{align*}
第一步用三角不等式，第二步用 $a^\dagger a=N_{\rm ph}$、$aa^\dagger=N_{\rm ph}+I$ 及 $0\le P_W\le I$。在
\[
 (U-U_0)\Psi_i=-\frac\ii\hbar\int U(t_f,t)R(t)\Psi_0(t)\dd t
\]
中插入 $P_W+(I-P_W)$，再用 $U$ 的幺正性，便得到\eqnrefs{eq:feqo-quantum-ladder-propagation-error}。
\end{proof}

因此相同的逐边失谐在高光子占据下会产生更大的振幅误差；只检查电子反冲频率和作用时间不足以给出所有光态的一致精度。若参考数矩至多为 $\mathcal N$，相对耦合差至多为 $\epsilon_G$，失谐差至多为 $d_\Delta$，则前一项至多为 $2\sqrt{\mathcal N+1}\int|\mathcal G_0(t)|(\epsilon_G+\tau d_\Delta)\dd t$。强脉冲中的累计耦合、玻色增强和工作窗尾部由此都进入同一个可核算的误差界。
